\section{Introduction}\label{sec:introduction}

For a smooth projective complex variety \(X\), the rational Hodge
conjecture asks whether every rational cohomology class of type
\((p,p)\) is a rational linear combination of codimension-\(p\)
algebraic cycle classes.  We prove this assertion for abelian
varieties with complex multiplication, in every dimension and
codimension.

A complex abelian variety \(A\) is of \emph{CM type}, or simply
\emph{CM}, if \(\End(A)\otimes_{\Z}\Q\) contains a commutative
semisimple \(\Q\)-algebra of dimension \(2\dim A\).  We include the
point by convention.  Finite products and powers of CM abelian
varieties are again CM.

\Needspace{12\baselineskip}
\begin{theorem}\label{thm:cm}
Let \(A\) be a complex CM abelian variety.  For every integer
\(p\geq0\), the Betti cycle-class map
\[
 \operatorname{cl}_B:\CH^p(A)_{\Q}
 \longrightarrow H^{2p}(A,\Q)\cap H^{p,p}(A)
\]
is surjective, where \(\CH^p(A)_{\Q}=\CH^p(A)\otimes_{\Z}\Q\) and
the intersection is taken in \(H^{2p}(A,\C)\).  In particular, the
rational Hodge conjecture holds in every codimension on every finite
product of complex CM abelian varieties and on every power of such a
product.
\end{theorem}

The theorem also has consequences beyond the algebraicity of individual
CM Hodge classes.  The generalized Hodge conjecture predicts algebraic
supports of the codimension allowed by the Hodge types of a cohomology
substructure.  We obtain it for CM abelian varieties by applying the
theorem to correspondences from auxiliary CM varieties, following
Hazama and Abdulali \cite{Hazama2003CM,Abdulali2005CM}.
Two theorems of Milne then give the Tate conjecture for all abelian
varieties over finite fields, which asks for algebraic representatives
of Galois-invariant \(\ell\)-adic classes, and the Hodge standard
conjecture for abelian varieties in arbitrary characteristic, which
asserts positivity of the intersection form on primitive algebraic
classes \cite{Milne1999Tate,Milne2002Polarizations}.
Section~\ref{cc:section} states these conclusions with their coefficient
and support conventions and explains the separate implications.

\subsection{Hodge classes on CM abelian varieties}

The codimension-one case of the Hodge conjecture is the Lefschetz
\((1,1)\) theorem.  Hodge formulated the broader problem by asking
for an extension of this algebraicity criterion to higher-dimensional
cycles \cite[p.~184]{Hodge1952}.  Cup products of divisor classes therefore give
algebraic Hodge classes in higher degree.  For an abelian variety,
\(H^*(A,\Q)=\Lambda^*H^1(A,\Q)\), but this description of its
cohomology does not make every Hodge class a linear combination of
divisor products.  The problem is to construct cycles for the remaining
classes.

Weil classes provide an important instance of this difficulty.  Suppose
an imaginary quadratic field \(K\) acts on an abelian variety of
dimension \(2n\), with its two embeddings occurring with equal
multiplicity in \(H^{1,0}\).  Then
\(\Lambda_K^{2n}H^1(A,\Q)\) supplies a two-dimensional rational
space of classes of type \((n,n)\).  These determinant classes and
their generalization to CM fields are discussed in
\cite[Proposition~4.4]{Deligne1982}.  Their Hodge type follows from
the multiplicity condition; algebraicity requires a geometric
construction.  Weil classes need not lie in the algebra generated by
divisor classes; this occurs already on abelian fourfolds
\cite[Theorem~0.1]{MoonenZarhin1999}.  Recent work of Markman proves
algebraicity of the Weil classes on polarized abelian sixfolds of
Weil type for an imaginary quadratic field, with polarization
discriminant \(-1\), and deduces the Hodge conjecture for all abelian
fourfolds \cite[Theorem~1.5.1 and Corollary~1.6.1]{Markman2025Secant}.
The sixfold theorem assumes an imaginary quadratic action, without
requiring the full CM condition of Theorem~\ref{thm:cm}.

Deligne proved that every Hodge class on an abelian variety is
\emph{absolute Hodge}: after transport to a conjugate variety through
de Rham cohomology, it remains a Hodge class
\cite[Main Theorem~2.11]{Deligne1982}.  This establishes an arithmetic
compatibility shared by algebraic cycle classes, but does not by itself
produce algebraic cycles.  CM varieties play a central role in that
argument.  Their Hodge groups are tori, and their degree-one cohomology
splits, after extending coefficients, into lines indexed by field
embeddings.  An embedding-line tensor belongs to the scalar extension
of the rational Hodge classes exactly when every scalar conjugate
has equally many holomorphic and antiholomorphic factors.
This is Pohlmann's balance criterion \cite{Pohlmann1968}; see
\cite[\S1.2(c)]{Milne2022} and
\cite[Theorem~3.1]{GaoUllmo2026}.

The CM arguments of Deligne and Andr\'e reduce these Hodge classes to
Weil-type classes on auxiliary CM abelian varieties
\cite{Deligne1982,Andre1992}.  More precisely, the formulation in
\cite[Theorem~1]{Milne2022}, attributed there to Andr\'e, expresses a
CM Hodge class as a sum of pullbacks of split Weil classes.  Here
\emph{split} refers to the Hermitian polarization form having a totally
isotropic subspace of half its dimension.  This reduction identifies
classes whose algebraicity would suffice; it does not supply their
algebraic representatives in all dimensions.

There is also a reduction to relations involving only four labels.
Hazama's arrangement method spans its rational balance kernel by
antipodal pairs and four-label square relations
\cite[Theorem~7.13]{Hazama2003CM}; his Corollary~7.14.1 gives the
corresponding codimension-two algebraicity reduction.  Gao and Ullmo
give equivalent quadratic indicator relations and use them, after
adjoining auxiliary CM factors, to generate Hodge-induced period
relations \cite[Theorems~7.1--7.2]{GaoUllmo2026}.
Their conclusion concerns relations among periods.  Our task is to
construct algebraic representatives for the corresponding four-factor
tensors and to compose those representatives as correspondences.
This distinction matters: the reduction uses auxiliary factors and
polarization contractions, rather than asserting that low-codimension
classes generate the Hodge ring of a fixed power.

\subsection{From balanced tensors to surface periods}

Choose a Galois CM field \(E\) containing the CM fields in
\(H^1(A,\Q)\), and write \(G_E=\Gal(E/\Q)\) and \(c\) for its
complex conjugation.  An odd sign function
\(v:G_E\to\{1,-1\}\), satisfying \(v(c\tau)=-v(\tau)\), specifies
a weight-one Hodge structure \(U(v)\) on the rational vector space
\(E\).  Its embedding line \(U(v)_\lambda\), on which \(a\in E\)
acts as \(\lambda(a)\), has type \((1,0)\) or \((0,1)\) according
as \(v(\lambda)\) is positive or negative.
Section~\ref{sec:tensors} realizes \(U(v)=H^1(B(v),\Q)\) for a CM
abelian variety \(B(v)\).

The elementary relation needed in the proof is
\[
 v_1+v_2=v_3+v_4.
\]
It makes every scalar conjugate of the four-factor line
\[
 U(v_1)_1\otimes U(v_2)_1\otimes U(v_3)_c\otimes U(v_4)_c
 \ \subset H^4\!\left(\prod_{i=1}^4B(v_i),\Qbar\right)
\]
have type \((2,2)\).  We prove that this line lies in the scalar
extension of the rational cycle span.  Section~\ref{sec:tensors}
then assembles arbitrary balanced tensors from these relations:
two-row switches connect lists with the same sum, and algebraic
polarization pairings contract their intermediate factors.  Transfer
back to the original abelian variety and descent of coefficients
give Theorem~\ref{thm:cm}.

The geometric construction of a four-factor class takes place on a
surface.  Suppose there are rational Hodge maps
\(h_i:U(v_i)\to H^1(S,\Q)\), for a smooth projective surface
\(S\), and vectors \(\alpha_i\) on the four selected embedding
lines such that
\[
 \int_S h_1(\alpha_1)\smile h_2(\alpha_2)
              \smile h_3(\alpha_3)\smile h_4(\alpha_4)\ne0.
\]
Degree-one Hodge maps are induced, up to rational multiples, by
morphisms \(S\to B(v_i)\).  The image of \(S\) in their product
therefore gives an algebraic cycle whose pairing with the desired
line is nonzero.  Section~\ref{sec:surface} shows how four algebraic
divisor kernels convert this functional into a vector on that exact
line.  The remaining problem is thus to find four degree-one classes
with both a nonzero surface period and the prescribed rational CM
sources.

\subsection{Theta forms and their CM sources}

Section~\ref{sec:theta} constructs the period on a compact arithmetic
quotient of the complex two-ball.  Its unitary group has signature
\((2,1)\) at one distinguished embedding and is definite at all the
other real places.  After the elementary normalization in
Section~\ref{sec:tensors}, all four lines \(U(v_i)_1\) are holomorphic
at that embedding.  To each type the construction assigns a family
of holomorphic theta one-forms.  Their values span a cotangent fiber,
so two such families supply a wedge that is a nonzero holomorphic
two-form.

To compare the two sides of the four-factor relation, we realize each
pair of types by an orthogonal decomposition of the same Hermitian
plane.  The balance relation says that the two decompositions have
the same local sign multiplicities.  Weak approximation produces a
rational second decomposition whose compact-place data are as close
as desired to those of the first.  At the distinguished embedding,
the total period data depend only on the plane and remain unchanged.

A rank-two theta tensor packages the wedges into one total two-form.
We keep its finite theta input and the compact quotient fixed while
varying the compact-place data.  The total form varies continuously,
so its pairing with the initial nonzero form stays close to the
nonzero self-pairing.  At the two rational endpoints it is a finite
sum of wedges of the required one-forms.  Factoring these endpoints
and passing to a common finite cover therefore gives a nonzero mixed
period with two holomorphic and two antiholomorphic factors.

The construction uses Weil's rational theta covariance and the unitary
splittings of the Weil representation
\cite{Weil1964,Kudla1994,BoradeEtAl2025ToriPeriods}.
Kudla and Millson developed differential-form theta kernels in this
unitary setting \cite[Introduction and Theorem~3]{KudlaMillson1990DifferentialTheta}.
The projection turning the rank-two tensor into wedges has a close
predecessor in the pure-cocycle product calculation of Bergeron,
Millson, and M\oe glin \cite[Lemmas~5.7 and~5.11]{BMM2016}.
We derive the required tensor identity in coordinates and use the
fixed-input variation to obtain the mixed period.

These are initially complex one-form classes.  The surface criterion
also requires that they come from specified rational CM Hodge
structures.  Sections~\ref{sec:arithmetic} and~\ref{sec:cm-sources}
establish this arithmetic input.  First we realize the ball quotients
in a projective moduli space and construct algebraic Hecke
endomorphisms of the product of its components' Albanese varieties.
Counting ordinary isogeny
neighbors gives a polynomial relation between these operators and
Frobenius.  We then compute the Hecke parameters of the theta forms.
The Frobenius relation and the degree-one Weil weight force slope
zero or one on the relevant Hecke summands, with the choice determined
by the prescribed CM signs.

The Hecke projectors preserve the Hodge filtration, so weak
admissibility turns these extreme slopes into pure Hodge types.
Applying the calculation
to every scalar conjugate identifies the complete CM type of each
summand.  Scalar descent then expresses every relevant theta class
as a complex linear combination of images of the selected CM line
under rational Hodge maps.  Section~\ref{sec:assembly} puts the
finite families on one projective surface, expands the nonzero
period, and applies the surface criterion.  This completes the
four-factor construction and hence the proof of the main theorem.
